\scanpage{43}
\mspage{34}

A-t-on \unsure{même}, pour $Y,Z\in\widetilde X$
\[
|Y|\subset|Z|\Longrightarrow Y\leq Z\quad ??
\]
Ça ne semble pas contenu dans (i) (ii) (iii) (iv), \deleted{À moins. Cela l'est.} Est-il impliqué par (ii bis)\unsure{??}
\deleted{J'explicite} À vrai dire, je l'ai sous-entendu en y disant ``indexé par $\widetilde X$'' \deleted{\unsure{comme}} comme \uline{un \unsure{ordonné}}, je présume) --- mais peut-on le prouver ? \deleted{Soit $L=\widetilde Y\cap\widetilde Z$} Ce qui est sûr, c'est que j'aurai besoin de savoir que pour $Y\in\widetilde X$,
\[
\begin{aligned}
|Y|^{\circ}\ &\text{est aussi \unsure{l'ens. analogue}, défini}\\
&\text{en termes de la figure ens.}\\
&\text{indexée par }\widetilde X,\text{ i.e. que l'on a}
\end{aligned}
\]
\[
\overset{\inserted{\unsure{\text{les}|\text{oui}}}}{|Y|}^{\circ}\overset{?}{=}\overset{\inserted{\unsure{\text{les}|\text{oui}}}}{|Y|}\setminus
\bigcup_{\substack{Z\in\widetilde X\ \text{tel que}\\ |Z|\unsure{\not\supset|\not\subset}|Y|\\\deleted{\unsure{\text{aussi}}}}}|Z|
\overset{\mathrm{def}}{=} |Y|^{\circ\prime}
\]
mais dans la $\bigcup$, il y a a priori \inserted{\uline{peut-être}} plus de $Z$ que \deleted{\unsure{avec}} les seuls $Z<Y$, donc le \deleted{\unsure{2}} deuxième \unsure{mb} risque d'être \uline{plus petit} que le premier.. Mais (ii bis) \inserted{(forme faible)} nous affirme que les $|Y|^{\circ\prime}$ recouvrent déjà $|X|$ (et sont tous non vides). A fortiori les $|Y|^{\circ}$ recouvrent \deleted{donc} Mais on sait qu'ils sont deux à deux disjoints par (iii) --- d'où finalement l'égalité $|Y|^{\circ}=|Y|^{\circ\prime}$, et la relation
\[
|X|=\bigcup_{Y\in\widetilde X}|Y|^{\circ}.
\]

Revenons alors à la question si
\[
|Y|\subset|Z|\Longrightarrow Y\leq Z,
\]
introduisons
\[
\widetilde Y\cap\widetilde Z=L,\ \deleted{\text{on a donc}}\ \widetilde Y\cup\widetilde Z=M
\]
% archive sheet number: 43
